\documentclass[10pt,a4paper]{amsart}
\usepackage[margin=19mm]{geometry}
\usepackage{amsmath,amssymb,mathtools}
\usepackage[scr=rsfs,cal=euler]{mathalfa}
\usepackage{xfrac}
\usepackage[unicode,hidelinks,bookmarks=false]{hyperref}
\newcommand{\ie}{i.e.\ }
\newcommand{\cf}{cf.\ }
\newcommand{\resp}{respectively\ }
\newcommand{\Subsection}{\subsection}
\providecommand{\nobreakdash}{\nobreak\hbox{-}}
\setlength{\emergencystretch}{2em}
\multlinegap0pt
\usepackage{bm}
\usepackage{braket}
\usepackage{amssymb}
\usepackage{mathtools}
\usepackage[shortlabels]{enumitem}
\usepackage{xcolor}
\usepackage{tikz}
\newtheorem{lemma}{Lemma}
\newcommand{\CC}{\ensuremath{\mathbb{C}}}
\newcommand{\NN}{\ensuremath{\mathbb{N}}}
\newcommand{\bfa}{\ensuremath{\mathbf{a}}}
\newcommand{\bfb}{\ensuremath{\mathbf{b}}}
\newcommand{\bfk}{\ensuremath{\mathbf{k}}}
\newcommand{\bfx}{\ensuremath{\mathbf{x}}}
\title{Unobstructedness of Affine Terminal Gorenstein Toric Fourfolds}
\author{Matej Filip, Alja\v{z} Zalar}
\date{Source: arXiv:2608.24424v1 (CC BY 4.0)}
\begin{document}
\maketitle

\noindent\emph{Source and licence.} Unobstructedness of Affine Terminal Gorenstein Toric Fourfolds. Version arXiv:2608.24424v1 (CC BY 4.0), distributed by arXiv under the Creative Commons Attribution 4.0 International licence, http://creativecommons.org/licenses/by/4.0/, as recorded in the arXiv licence metadata for this exact version and shown on the abstract page https://arxiv.org/abs/2608.24424v1. Full licence notice: \texttt{licenses/arxiv-2608.24424-CC-BY-4.0.txt}.\par\smallskip

\noindent\textit{Editorial scope.} Counted unit: the article's lemma boundary (source0.tex lines 2060--2069). The article's complete original proof of it is delivered with it (source0.tex lines 2071--2121, one \texttt{proof} environment of 51 lines). No mathematics is written, repaired or re-derived: the statement and the proof are the article's own lines, in the article's own notation, and no retained line is substituted or re-flowed.  Retained uncounted as the source-local context the proof needs: lines 1503--1551; lines 1844--1857.  Nothing was omitted from any retained span, so the package carries no omission record.  No retained line contains a citation, so the wrapper carries no bibliography and no bibliography entry.  No retained line contains a figure environment, an image-inclusion command or drawing code, so no image is delivered and the package declares no asset dependency.  Editorial formatting only, in addition: an AMS \texttt{amsart} preamble that restates the article's own declarations and macros used by the retained lines, namely the environments \texttt{lemma} on separate counters, together with the standard packages those lines need.  The retained environments print in this document as Lemma 1 in order of appearance, whereas the article prints the same items with the numbering of its own full document.  The complete native source of the article as distributed in the arXiv e-print for this exact version is bundled unchanged as \texttt{sources/208474/source0.tex}.
\begin{lemma}\label{lem presentation arbitrary generators}
For any finite set of boundary generators as above and any choices of
boundary representatives \(\partial(b)\), the binomials
\[
f_{\bfk}
:=
\bfx^{\bfk}-\chi_P^{s_{\bfk}},
\qquad
\bfk\in\NN^r,
\]
generate \(I_P\). 
{Let
$
\mathcal P:=\CC[\bfx,u]
$
and let
\[
\mathcal F
:=
\bigoplus_{\bfk\in\NN^r}\mathcal P e_{\bfk}
\]
be the free \(\mathcal P\)-module with basis
\(\{e_{\bfk}\mid\bfk\in\NN^r\}\). Let
\[
\psi\colon\mathcal F\longrightarrow I_P
\]
be the unique \(\mathcal P\)-linear homomorphism determined by
\[
\psi(e_{\bfk})=f_{\bfk}.
\]
Then \(\psi\) is surjective, and its kernel is generated as a
\(\mathcal P\)-module by the elements
\[
\rho_{\bfa,\bfk}
:=
e_{\bfa+\bfk}
-
\bfx^{\bfa}e_{\bfk}
-
u^{\eta_P(\bfk)}
e_{\partial(\bfk)+\bfa},
\qquad
\bfa,\bfk\in\NN^r.
\]
Equivalently, the first syzygy module of the family
\(\{f_{\bfk}\}_{\bfk\in\NN^r}\) is generated by the
\(\rho_{\bfa,\bfk}\).
}
\end{lemma}

\begin{lemma}\label{lem adapted boundary representatives}
After adjoining finitely many boundary elements to
\(s_1,\ldots,s_r\), every boundary element
\(b\in\partial_P(S_P)\) admits a representative \(\partial(b)\) such
that
\begin{equation}\label{eq adapted boundary vanishing}
\eta_{Q_i}(\partial(b))=0,
\qquad i=1,2.
\end{equation}
The enlarged elements together with \(R^*\) still generate \(S_P\),
and the presentation and syzygy statement of
Lemma~\ref{lem presentation arbitrary generators} applies to this
enlarged presentation.
\end{lemma}

\begin{lemma}
\label{lem eta compatible boundary}
For every \(\bfa\in\NN^r\),
\[
\eta_{Q_i}(\bfk_{\boldsymbol\alpha}+\bfa)
=
\eta_{Q_i}(\partial(\bfk)+\bfa),
\qquad i=1,2.
\]
\end{lemma}

\begin{proof}
Fix \(i\in\{1,2\}\), and write
\[
c_{\bfb}:=\sum_{j=1}^r b_jc_j
\qquad
\text{for }\bfb\in\NN^r.
\]
Since \(\partial(\bfk)\) and \(\bfk_{\boldsymbol\alpha}\) are the chosen
adapted representatives of boundary elements,
Lemma~\ref{lem adapted boundary representatives} gives
\[
\eta_{Q_i}(\partial(\bfk))
=
\eta_{Q_i}(\bfk_{\boldsymbol\alpha})
=
0.
\]
Moreover, projecting the boundary decomposition to \(M\) yields
$
c_{\bfk_{\boldsymbol\alpha}}
=
c_{\partial(\bfk)}
-
\pi_M(m_{\boldsymbol\alpha}).
$
By mutual compatibility,
\[
Q_i\subset\ker\pi_M(m_1)\cap\ker\pi_M(m_2),
\]
and hence \(\pi_M(m_{\boldsymbol\alpha})\) vanishes identically on
\(Q_i\). Therefore
$
\eta_{Q_i}(c_{\bfk_{\boldsymbol\alpha}})
=
\eta_{Q_i}(c_{\partial(\bfk)})
$
and
\[
\eta_{Q_i}(c_{\bfa}+c_{\bfk_{\boldsymbol\alpha}})
=
\eta_{Q_i}(c_{\bfa}+c_{\partial(\bfk)}).
\]
Using the definition of \(\eta_{Q_i}(\,\cdot\,)\) and the two vanishing
equalities above, we conclude that
\[
\eta_{Q_i}(\bfk_{\boldsymbol\alpha}+\bfa)
=
\eta_{Q_i}(\partial(\bfk)+\bfa),
\]
as required.
\end{proof}

\end{document}
